% GENERATED FILE. Edit canonical lessons, then run npm run generate:books.
% canonical-source-hash: fnv1a64:6c32623d1c62fbf9
% canonical-lessons: ES-C644-motor, ES-C644-volante, ES-C644-maletero, ES-C644-cinturon, ES-C644-abrocharse

\chapter{Dentro del coche}
\label{ch:es-dentro-del-coche}
\hlchaptermodality{644}

\begin{chapteropening}
\textbf{By the end of this chapter:} \emph{I can name the parts of a car I use on a trip: the engine, the steering wheel, the boot and the seat belt.}

Complete the last lesson of chapter 644: I can name the parts of a car I use on a trip: the engine, the steering wheel, the boot and the seat belt.
\end{chapteropening}

\section[el motor]{el motor — the engine, the motor}

\label{lesson:ES-C644-motor}



\begin{quote}
Before the new one: say the Spanish for to speed up, then the Spanish for to circulate.
\end{quote}

\subsection*{You'll want to know: el motor}
\textbf{el motor} — \textquotedblleft{}the engine, the motor\textquotedblright{}.

\begin{quote}
El motor hace un ruido raro.
\end{quote}

\begin{quote}
\emph{The engine is making a strange noise.}
\end{quote}

\emph{Arrancar el motor} — to start the engine.

\begin{cousinweb}
\emph{motor} is a learned borrowing from Latin \emph{mōtor} \textquotedblleft{}a mover, one that sets a thing moving\textquotedblright{} (Lewis \& Short; EtymDB), from \emph{movēre} \textquotedblleft{}to move\textquotedblright{}.
\end{cousinweb}

\begin{grammarlens}[title={What you've built}]
One more word for telling, explaining and describing at B1.
\end{grammarlens}

\subsection*{Your turn}
\begin{itemize}
\raggedright
  \item \emph{Say it:} \emph{el motor}
  \item \emph{Say it:} \emph{el motor}, once more
  \item \emph{Say it:} a sentence of your own with \emph{el motor}
  \item \emph{Recall:} say the Spanish for to speed up, then the Spanish for to circulate, then say \emph{el motor} again
\end{itemize}

\begin{tcolorbox}[breakable,colback=teal!4,colframe=teal!35!black,title={Before you move on}]
What does el motor mean? (\textquotedblleft{}The engine, the motor\textquotedblright{}.) Say it once more.
\end{tcolorbox}
\section[el volante]{el volante — the steering wheel}

\label{lesson:ES-C644-volante}



\begin{quote}
Before the new one: say the Spanish for to circulate, then the Spanish for the engine.
\end{quote}

\subsection*{You'll want to know: el volante}
\textbf{el volante} — \textquotedblleft{}the steering wheel\textquotedblright{}.

\begin{quote}
Pon las dos manos en el volante.
\end{quote}

\begin{quote}
\emph{Put both hands on the wheel.}
\end{quote}

\emph{Estar al volante} — to be at the wheel, to be driving.

\begin{cousinweb}
\emph{volante} is built on \emph{volar} \textquotedblleft{}to fly\textquotedblright{}, inherited from Latin \emph{volāre}; as an adjective it means \textquotedblleft{}flying\textquotedblright{} (EtymDB), and it came to name parts that turn or move freely.
\end{cousinweb}

\begin{grammarlens}[title={What you've built}]
One more word for telling, explaining and describing at B1.
\end{grammarlens}

\subsection*{Your turn}
\begin{itemize}
\raggedright
  \item \emph{Say it:} \emph{el volante}
  \item \emph{Say it:} \emph{el volante}, once more
  \item \emph{Say it:} a sentence of your own with \emph{el volante}
  \item \emph{Recall:} say the Spanish for to circulate, then the Spanish for the engine, then say \emph{el volante} again
\end{itemize}

\begin{tcolorbox}[breakable,colback=teal!4,colframe=teal!35!black,title={Before you move on}]
What does el volante mean? (\textquotedblleft{}The steering wheel\textquotedblright{}.) Say it once more.
\end{tcolorbox}
\section[el maletero]{el maletero — the boot (of a car), the trunk}

\label{lesson:ES-C644-maletero}



\begin{quote}
Before the new one: say the Spanish for the engine, then the Spanish for the steering wheel.
\end{quote}

\subsection*{You'll want to know: el maletero}
\textbf{el maletero} — \textquotedblleft{}the boot (of a car), the trunk\textquotedblright{}.

\begin{quote}
Las maletas van en el maletero.
\end{quote}

\begin{quote}
\emph{The suitcases go in the boot.}
\end{quote}

Spain says \emph{el maletero}; Mexico \emph{la cajuela}; much of South America \emph{el baúl}.

\begin{cousinweb}
\emph{maletero} is built on \emph{maleta} \textquotedblleft{}suitcase\textquotedblright{}. \emph{Maleta} is usually traced to French \emph{mallette}, a little \emph{malle} \textquotedblleft{}trunk\textquotedblright{}, which goes back to a Frankish word for a bag.
\end{cousinweb}

\begin{grammarlens}[title={What you've built}]
One more word for telling, explaining and describing at B1.
\end{grammarlens}

\subsection*{Your turn}
\begin{itemize}
\raggedright
  \item \emph{Say it:} \emph{el maletero}
  \item \emph{Say it:} \emph{el maletero}, once more
  \item \emph{Say it:} a sentence of your own with \emph{el maletero}
  \item \emph{Recall:} say the Spanish for the engine, then the Spanish for the steering wheel, then say \emph{el maletero} again
\end{itemize}

\begin{tcolorbox}[breakable,colback=teal!4,colframe=teal!35!black,title={Before you move on}]
What does el maletero mean? (\textquotedblleft{}The boot (of a car), the trunk\textquotedblright{}.) Say it once more.
\end{tcolorbox}
\section[el cinturón]{el cinturón — the belt; the seat belt}

\label{lesson:ES-C644-cinturon}



\begin{quote}
Before the new one: say the Spanish for the steering wheel, then the Spanish for the boot.
\end{quote}

\subsection*{You'll want to know: el cinturón}
\textbf{el cinturón} — \textquotedblleft{}the belt; the seat belt\textquotedblright{}.

\begin{quote}
Ponte el cinturón, por favor.
\end{quote}

\begin{quote}
\emph{Put your seat belt on, please.}
\end{quote}

\emph{El cinturón de seguridad} — the seat belt.

\begin{cousinweb}
\emph{cinturón} is built on \emph{cintura} \textquotedblleft{}waist\textquotedblright{}, which is inherited from Latin \emph{cinctūra} \textquotedblleft{}a girdle\textquotedblright{}, from \emph{cingere} \textquotedblleft{}to put a belt round\textquotedblright{}.
\end{cousinweb}

\begin{grammarlens}[title={What you've built}]
One more word for telling, explaining and describing at B1.
\end{grammarlens}

\subsection*{Your turn}
\begin{itemize}
\raggedright
  \item \emph{Say it:} \emph{el cinturón}
  \item \emph{Say it:} \emph{el cinturón}, once more
  \item \emph{Say it:} a sentence of your own with \emph{el cinturón}
  \item \emph{Recall:} say the Spanish for the steering wheel, then the Spanish for the boot, then say \emph{el cinturón} again
\end{itemize}

\begin{tcolorbox}[breakable,colback=teal!4,colframe=teal!35!black,title={Before you move on}]
What does el cinturón mean? (\textquotedblleft{}The belt; the seat belt\textquotedblright{}.) Say it once more.
\end{tcolorbox}
\section[abrocharse]{abrocharse — to fasten (a belt, a button)}

\label{lesson:ES-C644-abrocharse}



\begin{quote}
Before the new one: say the Spanish for the boot, then the Spanish for the belt.
\end{quote}

\subsection*{You'll want to know: abrocharse}
\textbf{abrocharse} — \textquotedblleft{}to fasten (a belt, a button)\textquotedblright{}.

\begin{quote}
Abróchense el cinturón, por favor.
\end{quote}

\begin{quote}
\emph{Please fasten your seat belts.}
\end{quote}

Without \emph{se}, \emph{abrochar} is to fasten something: \emph{abrochar el abrigo}.

\begin{cousinweb}
\emph{abrocharse} is built on \emph{broche} \textquotedblleft{}clasp, brooch\textquotedblright{}, which was borrowed from French \emph{broche} (EtymDB).
\end{cousinweb}

\begin{grammarlens}[title={What you've built}]
One more word for telling, explaining and describing at B1.
\end{grammarlens}

\subsection*{Your turn}
\begin{itemize}
\raggedright
  \item \emph{Say it:} \emph{abrocharse}
  \item \emph{Say it:} \emph{abrocharse}, once more
  \item \emph{Say it:} a sentence of your own with \emph{abrocharse}
  \item \emph{Recall:} say the Spanish for the boot, then the Spanish for the belt, then say \emph{abrocharse} again
\end{itemize}

\begin{tcolorbox}[breakable,colback=teal!4,colframe=teal!35!black,title={Before you move on}]
What does abrocharse mean? (\textquotedblleft{}To fasten (a belt, a button)\textquotedblright{}.) Say it once more.
\end{tcolorbox}
